% Stat 221 LaTeX preamble (shared across lecture notes)

% --- Page + typography ---
\usepackage[margin=1in]{geometry}
\usepackage[T1]{fontenc}
\usepackage{libertinus}
\usepackage{libertinust1math}
\usepackage[scaled=0.95]{inconsolata}
\usepackage{microtype}

% --- Math + symbols ---
\usepackage{amsmath,amssymb,amsthm,mathtools}
\usepackage{bm}

% --- Layout helpers ---
\usepackage{graphicx}
\usepackage{booktabs}
\usepackage{enumitem}
\usepackage{titlesec}
\usepackage{fancyhdr}
\usepackage{xcolor}
\usepackage[most]{tcolorbox}
\usepackage{algorithm}
\usepackage{algpseudocode}

% --- Visual style ---
\definecolor{stat221Blue}{HTML}{1F4E79}
\definecolor{stat221BlueLight}{HTML}{EAF3FA}
\definecolor{stat221GrayLight}{HTML}{F6F7FB}
\definecolor{stat221Gray}{HTML}{2F2F2F}
\definecolor{stat221Purple}{HTML}{6A3D9A}
\definecolor{stat221PurpleLight}{HTML}{F3EEFA}
\definecolor{stat221Gold}{HTML}{9A6B00}
\definecolor{stat221GoldLight}{HTML}{FFF3D9}
\definecolor{stat221Teal}{HTML}{0F766E}
\definecolor{stat221TealLight}{HTML}{EAF7F6}
\definecolor{stat221Green}{HTML}{2E7D32}
\definecolor{stat221GreenLight}{HTML}{EDF8EF}
\definecolor{stat221Orange}{HTML}{B45309}
\definecolor{stat221OrangeLight}{HTML}{FFF4E8}
\definecolor{stat221Red}{HTML}{B2182B}
\colorlet{stat221TitleBack}{stat221Blue!12}

% --- Figures ---
\usepackage{tikz}
\usetikzlibrary{arrows.meta,calc,positioning,matrix}
\usepackage{pgfplots}
\pgfplotsset{compat=1.18}

% --- Bibliography ---
\usepackage[round]{natbib}

% --- Links and cross-references ---
\usepackage[colorlinks=true,linkcolor=stat221Blue,citecolor=stat221Blue,urlcolor=stat221Blue]{hyperref}
\usepackage[nameinlink,capitalise,noabbrev]{cleveref}

% Better names in references.
\crefname{algorithm}{Algorithm}{Algorithms}
\Crefname{algorithm}{Algorithm}{Algorithms}

% Refer to all sectioning levels uniformly as "Section".
\crefname{section}{Section}{Sections}
\Crefname{section}{Section}{Sections}
\crefname{subsection}{Section}{Sections}
\Crefname{subsection}{Section}{Sections}
\crefname{subsubsection}{Section}{Sections}
\Crefname{subsubsection}{Section}{Sections}

% Algorithmicx wording.
\algrenewcommand\algorithmicrequire{\textbf{Input:}}
\algrenewcommand\algorithmicensure{\textbf{Output:}}

% Avoid duplicate hyperref destinations for algorithmic line numbers.
\makeatletter
\providecommand{\theHALG@line}{}
\renewcommand{\theHALG@line}{ALG@line.\thealgorithm.\arabic{ALG@line}}
\makeatother

\setlength{\parindent}{0pt}
\setlength{\parskip}{0.5em}

\setlist[itemize]{leftmargin=*,topsep=0.25em,itemsep=0.15em}
\setlist[enumerate]{leftmargin=*,topsep=0.25em,itemsep=0.15em}

\titleformat{\section}{\Large\bfseries\color{stat221Blue}}{\thesection}{0.75em}{}
\titleformat{\subsection}{\large\bfseries\color{stat221Blue}}{\thesubsection}{0.75em}{}
\titleformat{\subsubsection}{\bfseries\color{stat221Gray}}{\thesubsubsection}{0.75em}{}
\titleformat{\part}[display]
  {\centering\Huge\bfseries\color{stat221Blue}}
  {}
  {0pt}
  {\vspace*{\fill}\partname~\thepart\\[0.6em]}
  [\vspace*{\fill}\thispagestyle{plain}\clearpage]
\titlespacing*{\part}{0pt}{0pt}{0pt}
\titlespacing*{\section}{0pt}{2.2ex plus 0.4ex minus 0.2ex}{1.0ex plus 0.2ex}
\titlespacing*{\subsection}{0pt}{1.6ex plus 0.3ex minus 0.2ex}{0.6ex plus 0.2ex}
\titlespacing*{\subsubsection}{0pt}{1.2ex plus 0.2ex minus 0.2ex}{0.3ex plus 0.2ex}

% Make parts full-page dividers: ensure `\part{...}` always starts a new page.
\let\statOrigPart\part
\renewcommand{\part}{\clearpage\statOrigPart}

% Keep the document structure uniform:
% - Show sections, subsections, and subsubsections in the table of contents.
\setcounter{secnumdepth}{3}
\setcounter{tocdepth}{3}

\pagestyle{fancy}
\fancyhf{}
\lhead{\textsc{Stat 221}}
\rhead{\nouppercase{\leftmark}}
\cfoot{\thepage}
\renewcommand{\headrulewidth}{0.4pt}
\setlength{\headheight}{14pt}
\fancypagestyle{plain}{
  \fancyhf{}
  \cfoot{\thepage}
  \renewcommand{\headrulewidth}{0pt}
}

% --- Theorem environments ---
\theoremstyle{definition}
\newtheorem{definition}{Definition}[section]
\newtheorem{example}[definition]{Example}
\newtheorem{exercise}[definition]{Exercise}

\theoremstyle{plain}
\newtheorem{theorem}[definition]{Theorem}
\newtheorem{proposition}[definition]{Proposition}
\newtheorem{lemma}[definition]{Lemma}
\newtheorem{corollary}[definition]{Corollary}

\theoremstyle{remark}
\newtheorem{remark}[definition]{Remark}

% --- Boxed theorem styling (keeps amsthm counters) ---
\tcbset{
  stat221box/.style={
    enhanced,
    breakable,
    lines before break=3,
    sharp corners,
    boxrule=0.6pt,
    coltitle=stat221Gray,
    fonttitle=\bfseries,
    left=8pt,right=8pt,top=6pt,bottom=6pt,
  },
  stat221titlebox/.style={
    stat221box,
    colbacktitle=stat221TitleBack,
    boxed title style={colback=stat221TitleBack,colframe=stat221Blue!60!black},
  },
  stat221panelblue/.style={
    stat221titlebox,
    colback=stat221BlueLight,
    colframe=stat221Blue!60!black,
  },
  stat221panelgray/.style={
    stat221titlebox,
    colback=stat221GrayLight,
    colframe=stat221Blue!60!black,
  },
}

\tcolorboxenvironment{definition}{stat221box,colback=stat221BlueLight,colframe=stat221Blue!60!black}
\tcolorboxenvironment{example}{stat221box,colback=stat221GoldLight,colframe=stat221Gold!70!black}
\tcolorboxenvironment{exercise}{stat221box,colback=stat221OrangeLight,colframe=stat221Orange!70!black}
\tcolorboxenvironment{theorem}{stat221box,colback=stat221GrayLight,colframe=stat221Gray!70!black}
\tcolorboxenvironment{lemma}{stat221box,colback=stat221GreenLight,colframe=stat221Green!60!black}
\tcolorboxenvironment{proposition}{stat221box,colback=stat221PurpleLight,colframe=stat221Purple!60!black}
\tcolorboxenvironment{corollary}{stat221box,colback=stat221TealLight,colframe=stat221Teal!60!black}
\tcolorboxenvironment{remark}{stat221box,colback=white,colframe=stat221Gray!50!black}

\renewcommand{\qedsymbol}{$\blacksquare$}

% --- Macros ---
\newcommand{\R}{\mathbb{R}}
\newcommand{\C}{\mathbb{C}}
\newcommand{\E}{\mathbb{E}}
\newcommand{\1}{\mathbf{1}}

\DeclareMathOperator*{\argmin}{arg\,min}
\DeclareMathOperator*{\argmax}{arg\,max}

\DeclarePairedDelimiter\norm{\lVert}{\rVert}
\DeclarePairedDelimiter\ip{\langle}{\rangle}

% A consistent Hessian notation (to avoid ambiguity with other uses of \nabla^2).
\newcommand{\Hess}{\nabla^2}
